% ---------------------------------------------------------------
% macros.tex
% Shared notation for the Japanese and English versions of the report.
% Both versions MUST \input this file so that the two documents use
% literally the same symbols.  Do not define notation anywhere else.
%
% Naming policy: avoid every control sequence already provided by the
% `physics' package (\tr, \Tr, \dd, \dv, \pdv, \qty, \abs, \norm,
% \expval, \ev, \comm, \acomm, \pb, \order, ...).
% ---------------------------------------------------------------

% --- Unified description ----------------------------------------
% The report treats a classical probability vector and a density matrix on the
% same footing; \vstate denotes either, and \gen denotes whichever generator
% acts on it.  These two symbols carry the unification of the three papers.
\newcommand{\vstate}{\varrho}           % unified state: p(t) classically, rho(t) quantum mechanically
\newcommand{\gen}{\mathcal{G}}          % unified generator: R classically, \Lop quantum mechanically
\newcommand{\ssv}{\vstate_{\mathrm{ss}}}% unified stationary state
\newcommand{\ipr}[2]{\langle #1 , #2 \rangle} % duality pairing between left and right modes

% --- Generators of the dynamics ---------------------------------
\newcommand{\Lop}{\mathcal{L}}          % Lindbladian / Liouvillian superoperator
\newcommand{\Rmat}{R}                   % classical Markov rate matrix (notation of Lu-Raz)
\newcommand{\Wmat}{\mathsf{W}}          % classical Markov rate matrix (generic)
\newcommand{\jump}{L}                   % Lindblad jump operator

% --- Stationary states ------------------------------------------
\newcommand{\pss}{\pi}                  % classical stationary distribution
\newcommand{\rss}{\rho_{\mathrm{ss}}}   % quantum steady state

% --- Spectral decomposition -------------------------------------
\newcommand{\reig}[1]{r_{#1}}           % right eigenvector (eigenmode) of the generator
\newcommand{\leig}[1]{\ell_{#1}}        % left eigenvector, dual to \reig
\newcommand{\ovl}[1]{a_{#1}}            % overlap coefficient of the initial state with mode #1
\newcommand{\gap}{\lambda_{2}}          % spectral gap: slowest decaying nonzero mode

% --- Temperatures -----------------------------------------------
\newcommand{\betaeff}{\beta_{\mathrm{eff}}}   % effective inverse temperature
\newcommand{\betabath}{\beta_{\mathrm{bath}}} % inverse temperature of the bath
\newcommand{\betainit}{\beta_{0}}             % initial inverse temperature (state preparation)

% --- Keldysh Green's functions ----------------------------------
\newcommand{\Gret}{G^{R}}               % retarded Green's function
\newcommand{\Gadv}{G^{A}}               % advanced Green's function
\newcommand{\Gkel}{G^{K}}               % Keldysh Green's function
\newcommand{\Ggtr}{G^{>}}               % greater Green's function
\newcommand{\Glss}{G^{<}}               % lesser Green's function
\newcommand{\Sret}{\Sigma^{R}}          % retarded self-energy
\newcommand{\Sgtr}{\Sigma^{>}}          % greater self-energy

% --- Model parameters -------------------------------------------
\newcommand{\Jsyk}{J}                   % SYK four-fermion coupling
\newcommand{\Vsb}{V}                    % system-bath coupling
\newcommand{\hyb}{\Gamma}               % hybridisation function of the resonant level
\newcommand{\bandw}{\Lambda}            % bath bandwidth (Lorentzian half-width)
\newcommand{\edot}{\varepsilon_{d}}     % energy of the resonant level

% --- Distance measures ------------------------------------------
\newcommand{\Dbath}{D_{\mathrm{bath}}}  % distance from the asymptotic bath distribution
\newcommand{\Deq}{D_{\mathrm{eq}}}      % distance from the equilibrium state at the same \betaeff

% --- Miscellaneous ----------------------------------------------
% House style: the differential is upright (\dd from `physics'), while the
% imaginary unit and Euler's number stay italic.  The macros below exist so
% that this convention is applied uniformly and can be changed in one place.
\newcommand{\ii}{i}                     % imaginary unit (italic, per house style)
\newcommand{\ee}{e}                     % Euler's number (italic, per house style)
\newcommand{\Feq}{f_{\mathrm{eq}}}      % Fermi-Dirac distribution
\newcommand{\Feff}{f_{\mathrm{eff}}}    % nonequilibrium (effective) distribution function
